\documentclass[aspectratio=169]{beamer}

\usetheme{Madrid}
\usecolortheme{default}
\usefonttheme{professionalfonts}

\usepackage[utf8]{inputenc}
\usepackage[T1]{fontenc}
\usepackage{lmodern}
\usepackage{amsmath, amssymb, amsfonts}
\usepackage{bm}
\usepackage{booktabs}
\usepackage{graphicx}
\usepackage{hyperref}
\usepackage{xcolor}

\hypersetup{
	colorlinks=true,
	linkcolor=blue,
	urlcolor=blue
}

\title[PQC Hard Problem Orientation]{Orientation to Post-Quantum Cryptographic Hard Problems}
\author{Your Name}
\institute{Your Institution / Research Group}
\date{\today}

\begin{document}
	
	%------------------------------------------------
	\begin{frame}
		\titlepage
	\end{frame}
	
	%------------------------------------------------
	\begin{frame}{Outline}
		\tableofcontents
	\end{frame}
	
	%------------------------------------------------
	\section{Motivation}
	
	\begin{frame}{Why Post-Quantum Cryptography?}
		\begin{itemize}
			\item Classical public-key cryptography relies heavily on problems such as integer factorization and discrete logarithms.
			\item Large-scale quantum computers would threaten these assumptions through algorithms such as Shor's algorithm.
			\item Post-quantum cryptography (PQC) seeks cryptographic constructions secure against both classical and quantum adversaries.
			\item The central question is therefore:
			\begin{block}{Core Question}
				Which mathematical problems remain computationally intractable even in the presence of quantum computation?
			\end{block}
		\end{itemize}
	\end{frame}
	
	%------------------------------------------------
	\begin{frame}{Role of Hard Problems in Cryptography}
		A cryptosystem is only as strong as the computational assumption on which it is built.
		
		\vspace{0.5em}
		In the post-quantum setting, hard problems serve as:
		\begin{itemize}
			\item the \textbf{foundation of security proofs},
			\item the \textbf{source of asymptotic and concrete hardness},
			\item the \textbf{bridge} between abstract mathematics and practical protocol design.
		\end{itemize}
		
		\vspace{0.5em}
		Thus, understanding PQC requires understanding its underlying hard problems.
	\end{frame}
	
	%------------------------------------------------
	\section{What Is a Hard Problem?}
	
	\begin{frame}{Definition of a Hard Problem}
		Informally, a computational problem is considered \emph{hard} if no efficient algorithm is known for solving it on average or in the worst case under relevant parameter choices.
		
		\vspace{0.75em}
		In cryptography, one typically asks:
		\begin{itemize}
			\item Is the problem difficult for \textbf{classical} algorithms?
			\item Is the problem difficult for \textbf{quantum} algorithms?
			\item Is the hardness supported by \textbf{reductions}, \textbf{evidence}, or \textbf{long-standing study}?
			\item Does the problem admit \textbf{efficient implementation} in real systems?
		\end{itemize}
	\end{frame}
	
	%------------------------------------------------
	\begin{frame}{Desirable Features of a Cryptographic Assumption}
		A useful cryptographic hardness assumption should ideally have:
		\begin{enumerate}
			\item \textbf{Strong evidence of intractability},
			\item \textbf{Resistance to known classical and quantum attacks},
			\item \textbf{Efficient parameterization},
			\item \textbf{Reduction-based security arguments},
			\item \textbf{Robustness under implementation constraints}.
		\end{enumerate}
		
		\vspace{0.5em}
		In practice, there is often a trade-off between theoretical elegance and implementation efficiency.
	\end{frame}
	
	%------------------------------------------------
	\section{Main Families of PQC Hard Problems}
	
	\begin{frame}{Major Families of Post-Quantum Hard Problems}
		The principal mathematical families include:
		\begin{itemize}
			\item \textbf{Lattice-based problems}
			\begin{itemize}
				\item Shortest Vector Problem (SVP)
				\item Closest Vector Problem (CVP)
				\item Learning With Errors (LWE)
				\item Module-LWE / Ring-LWE
			\end{itemize}
			\item \textbf{Code-based problems}
			\begin{itemize}
				\item Syndrome Decoding
				\item General decoding problems for linear codes
			\end{itemize}
			\item \textbf{Multivariate problems}
			\begin{itemize}
				\item Solving systems of multivariate polynomial equations over finite fields
			\end{itemize}
			\item \textbf{Hash-based assumptions}
			\begin{itemize}
				\item Preimage resistance, second-preimage resistance, collision resistance
			\end{itemize}
			\item \textbf{Isogeny-based problems} (historically important, though currently more fragile in practice)
		\end{itemize}
	\end{frame}
	
	%------------------------------------------------
	\section{Lattice-Based Hard Problems}
	
	\begin{frame}{Lattices: Geometric Background}
		A lattice $\Lambda \subset \mathbb{R}^n$ is the set of all integer linear combinations of basis vectors:
		\[
		\Lambda = \left\{ \sum_{i=1}^m z_i \bm{b}_i : z_i \in \mathbb{Z} \right\}.
		\]
		
		\vspace{0.5em}
		Fundamental computational questions include:
		\begin{itemize}
			\item finding a short nonzero vector,
			\item finding a lattice vector close to a target point,
			\item solving noisy linear relations derived from lattice structure.
		\end{itemize}
	\end{frame}
	
	%------------------------------------------------
	\begin{frame}{Shortest Vector Problem (SVP)}
		\begin{block}{SVP}
			Given a lattice basis for $\Lambda$, find a nonzero vector $\bm{v} \in \Lambda$ minimizing $\|\bm{v}\|$.
		\end{block}
		
		\vspace{0.5em}
		Why it matters:
		\begin{itemize}
			\item It is a canonical worst-case lattice problem.
			\item Approximate variants of SVP are central in complexity theory and cryptography.
			\item Many security reductions relate average-case cryptographic assumptions to worst-case lattice problems.
		\end{itemize}
	\end{frame}
	
	%------------------------------------------------
	\begin{frame}{Learning With Errors (LWE)}
		\begin{block}{LWE Problem}
			Given samples of the form
			\[
			(\bm{a}_i, b_i = \langle \bm{a}_i, \bm{s} \rangle + e_i \bmod q),
			\]
			recover the secret $\bm{s}$ or distinguish such samples from uniform.
		\end{block}
		
		\vspace{0.5em}
		Here:
		\begin{itemize}
			\item $\bm{a}_i \in \mathbb{Z}_q^n$ is public,
			\item $\bm{s} \in \mathbb{Z}_q^n$ is secret,
			\item $e_i$ is a small error term.
		\end{itemize}
		
		\vspace{0.5em}
		LWE is one of the most important assumptions in modern post-quantum cryptography.
	\end{frame}
	
	%------------------------------------------------
	\begin{frame}{Why LWE Is Foundational}
		LWE is especially attractive because:
		\begin{itemize}
			\item it admits strong reduction-based evidence from worst-case lattice problems,
			\item it supports encryption, signatures, key exchange, FHE, IBE, and more,
			\item it has natural structured variants such as Ring-LWE and Module-LWE.
		\end{itemize}
		
		\vspace{0.5em}
		Many standardized PQC schemes are based on this family of assumptions.
	\end{frame}
	
	%------------------------------------------------
	\section{Code-Based Hard Problems}
	
	\begin{frame}{Code-Based Assumptions}
		Code-based cryptography is built on the difficulty of decoding a random linear code.
		
		\begin{block}{Syndrome Decoding Problem}
			Given a parity-check matrix $H$, a syndrome $\bm{s}$, and a weight bound $w$, find an error vector $\bm{e}$ such that
			\[
			H \bm{e}^T = \bm{s}^T
			\]
			and $\mathrm{wt}(\bm{e}) \leq w$.
		\end{block}
		
		\vspace{0.5em}
		This problem is believed to be hard even for quantum adversaries, though parameters must be chosen carefully.
	\end{frame}
	
	%------------------------------------------------
	\begin{frame}{Strengths and Limitations of Code-Based Cryptography}
		\textbf{Strengths}
		\begin{itemize}
			\item Long history of cryptanalytic study,
			\item Conservative security posture,
			\item Fast low-level operations in some settings.
		\end{itemize}
		
		\vspace{0.5em}
		\textbf{Limitations}
		\begin{itemize}
			\item Large public keys in many constructions,
			\item More limited flexibility compared with lattice-based designs for advanced primitives.
		\end{itemize}
	\end{frame}
	
	%------------------------------------------------
	\section{Multivariate and Hash-Based Problems}
	
	\begin{frame}{Multivariate Polynomial Problems}
		Multivariate cryptography relies on the difficulty of solving systems of nonlinear polynomial equations over finite fields.
		
		\vspace{0.5em}
		General form:
		\[
		f_1(x_1,\dots,x_n)=\cdots=f_m(x_1,\dots,x_n)=0.
		\]
		
		\vspace{0.5em}
		These problems are attractive because:
		\begin{itemize}
			\item evaluation can be very efficient,
			\item signing can be fast in carefully designed systems.
		\end{itemize}
		
		\vspace{0.5em}
		However, the design space is delicate and vulnerable to algebraic attacks if structure is not controlled.
	\end{frame}
	
	%------------------------------------------------
	\begin{frame}{Hash-Based Assumptions}
		Hash-based cryptography relies on standard assumptions such as:
		\begin{itemize}
			\item preimage resistance,
			\item second-preimage resistance,
			\item collision resistance.
		\end{itemize}
		
		\vspace{0.5em}
		Hash-based signatures are among the most conservative post-quantum constructions because their security reduces to the security of the underlying hash function.
		
		\vspace{0.5em}
		Their main trade-offs often concern:
		\begin{itemize}
			\item signature size,
			\item state management (for stateful variants),
			\item performance in constrained environments.
		\end{itemize}
	\end{frame}
	
	%------------------------------------------------
	\section{Security Perspective}
	
	\begin{frame}{What Makes a PQC Problem Credible?}
		The credibility of a post-quantum hardness assumption depends on:
		\begin{itemize}
			\item \textbf{Cryptanalytic maturity}: how long and how deeply the problem has been studied;
			\item \textbf{Reduction theory}: whether security can be related to better-understood problems;
			\item \textbf{Quantum attack landscape}: whether known quantum algorithms provide substantial speedups;
			\item \textbf{Parameter discipline}: whether implementations preserve the intended hardness level.
		\end{itemize}
	\end{frame}
	
	%------------------------------------------------
	\begin{frame}{Classical vs.\ Quantum Attack Models}
		In practice, one evaluates hard problems under both:
		\begin{enumerate}
			\item \textbf{Classical attacks}
			\begin{itemize}
				\item combinatorial methods,
				\item algebraic attacks,
				\item lattice reduction,
				\item information-set decoding,
				\item Gröbner-basis techniques.
			\end{itemize}
			\item \textbf{Quantum attacks}
			\begin{itemize}
				\item Shor-type exponential attacks where applicable,
				\item Grover-type quadratic speedups,
				\item quantum adaptations of classical search or algebraic methods.
			\end{itemize}
		\end{enumerate}
	\end{frame}
	
	%------------------------------------------------
	\section{From Hard Problems to Schemes}
	
	\begin{frame}{From Assumptions to Cryptographic Constructions}
		A hard problem becomes practically relevant only when it supports efficient schemes:
		\begin{center}
			\textbf{Hard Problem} $\longrightarrow$ \textbf{Assumption} $\longrightarrow$ \textbf{Construction} $\longrightarrow$ \textbf{Implementation}
		\end{center}
		
		\vspace{0.75em}
		Examples:
		\begin{itemize}
			\item LWE / Module-LWE $\rightarrow$ key encapsulation mechanisms,
			\item lattice trapdoors $\rightarrow$ digital signatures,
			\item hash assumptions $\rightarrow$ stateless and stateful signatures,
			\item decoding problems $\rightarrow$ code-based encryption/KEMs.
		\end{itemize}
	\end{frame}
	
	%------------------------------------------------
	\begin{frame}{Practical Evaluation Criteria}
		In real deployments, we compare candidate assumptions and schemes by:
		\begin{itemize}
			\item claimed security level,
			\item public key and ciphertext/signature size,
			\item runtime and memory cost,
			\item implementation simplicity,
			\item side-channel resilience,
			\item confidence in long-term hardness.
		\end{itemize}
	\end{frame}
	
	%------------------------------------------------
	\section{Research Orientation}
	
	\begin{frame}{How to Study PQC Hard Problems}
		A productive orientation path includes:
		\begin{enumerate}
			\item Learn the precise mathematical formulation of each problem.
			\item Understand known classical and quantum algorithms against it.
			\item Study reduction results and hardness evidence.
			\item Examine how the assumption is instantiated in concrete schemes.
			\item Track implementation and parameter-selection issues.
		\end{enumerate}
	\end{frame}
	
	%------------------------------------------------
	\begin{frame}{Suggested Research Questions}
		Possible guiding questions for a new student or researcher:
		\begin{itemize}
			\item What type of hardness evidence is available: worst-case, average-case, heuristic, or empirical?
			\item How does algebraic structure improve efficiency, and what security risk does it introduce?
			\item Which attacks dominate the concrete security estimates?
			\item How should one compare assumptions across different PQC families?
			\item What are the open problems in reduction theory, parameter setting, and implementation security?
		\end{itemize}
	\end{frame}
	
	%------------------------------------------------
	\section{Conclusion}
	
	\begin{frame}{Conclusion}
		\begin{itemize}
			\item Post-quantum cryptography is fundamentally an exercise in identifying and exploiting robust hard problems.
			\item Lattice, code-based, multivariate, and hash-based assumptions provide distinct trade-offs.
			\item The study of PQC hard problems requires both mathematical rigor and cryptographic realism.
			\item A clear understanding of these assumptions is essential for evaluating present and future post-quantum schemes.
		\end{itemize}
	\end{frame}
	
	%------------------------------------------------
	\begin{frame}{Thank You}
		\centering
		{\Large Thank you for your attention.}
		
		\vspace{1em}
		Questions and discussion are welcome.
	\end{frame}
	
\end{document}